% Image credits for the photographs and AI illustrations of Book 2
% (University Chemistry, Year 1). Machine-readable license records live in
% images/book2/CREDITS.md; keep the two files in sync. The Book 2 writer
% owns this file: add one \item per photograph (in an itemize, after the
% first paragraph) as photographs land.
\chapter*{Kredit Gambar}

Gambar-gambar dalam buku ini adalah karya asli. Foto-foto, jika dipakai,
dipakai dengan lisensi bebasnya masing-masing, dengan terima kasih kepada
para pembuatnya; tautan sumber dan URL lisensi lengkap untuk setiap foto
tercantum dalam \texttt{images/book2/CREDITS.md} di repositori buku ini.
Piktogram bahaya adalah piktogram Sistem Harmonisasi Global untuk
Klasifikasi dan Pelabelan Bahan Kimia, sebagaimana digambar untuk paket
\texttt{ghsystem} dari \TeX\ Live.

\bigskip
Di samping foto-foto dan gambar-gambar asli, beberapa gambar adalah
ilustrasi buatan AI, yang dihasilkan dengan pembangkit gambar OpenAI dari
perintah yang ditulis oleh para penyunting buku ini, dan diperiksa
ketepatan kimianya sebelum dimuat. Perintah lengkap untuk setiap gambar
yang dihasilkan dicatat dalam \texttt{images/book2/ai/PROMPTS.md} di
repositori.

\bigskip
\noindent Foto-foto, menurut bab:
\begin{itemize}\small
  \item Bab~1: Niels Bohr, sekitar 1922, AB Lagrelius \& Westphal, domain
        publik (Wikimedia Commons).
  \item Bab~2: Linus Pauling, 1962, Yayasan Nobel, domain publik
        (Wikimedia Commons).
  \item Bab~5: spesimen tembaga alami, oleh Flipin, OG, CC0 (Wikimedia
        Commons).
  \item Bab~6: kristal halit, fluorit, dan intan oktahedral, oleh James
        St. John, CC BY 2.0; grafit, oleh Zbynek Burival, CC BY-SA 4.0;
        kristal salju oleh Wilson A. Bentley (1902), domain publik
        (semuanya dari Wikimedia Commons).
  \item Bab~8: Svante Arrhenius, 1909, Meisenbach Riffarth \& Co., domain
        publik (Wikimedia Commons).
  \item Bab~11: endapan perak klorida, bromida, dan iodida, oleh
        Rrausch1974, CC BY-SA 3.0 (Wikimedia Commons).
  \item Bab~12: suspensi tembaga(II) hidroksida, oleh Alvy16, CC BY 4.0;
        larutan tembaga amina, oleh Chemicalinterest, domain publik
        (keduanya dari Wikimedia Commons).
  \item Bab~13: Walther Nernst, 1889, Smithsonian Libraries, domain publik
        (Wikimedia Commons).
  \item Bab~16: Jacobus Henricus van 't Hoff, foto yang terbit pada 1911,
        domain publik (Wikimedia Commons).
  \item Bab~21: Victor Grignard, domain publik (Wikimedia Commons).
  \item Bab~25: Vladimir Markovnikov, potret dari sebuah obituari yang
        terbit pada 1905, domain publik (Wikimedia Commons).
  \item Bab~26: kolam litium di Salar de Atacama, NASA Earth Observatory,
        domain publik; natrium dalam minyak mineral, oleh Justin Urgitis,
        CC BY-SA 2.5; warna nyala, oleh Hegelrast, CC BY-SA 4.0 (semuanya
        dari Wikimedia Commons).
  \item Bab~27: bauksit pisolitik, oleh James St. John, CC BY 2.0; boraks,
        oleh Aram Dulyan, domain publik; kuarsa, oleh Stephanie Clifford,
        CC BY 2.0; fosfor putih, Hi-Res Images of Chemical Elements,
        CC BY 3.0; Fritz Haber, Yayasan Nobel, domain publik (semuanya dari
        Wikimedia Commons).
  \item Bab~28: penambangan belerang di kawah Kawah Ijen, oleh S\'emhur,
        CC BY-SA 4.0; klorin dalam ampul, oleh W. Oelen, CC BY-SA 3.0; brom
        dalam ampul, oleh Jurii, CC BY 3.0; kristal iodin, oleh Dnn87,
        CC BY 3.0 (semuanya dari Wikimedia Commons).
  \item Bab~29: bangku pemanas (bangku Kofler), oleh HBR, CC BY-SA 3.0;
        refraktometer Abbe, oleh Foreade, CC BY-SA 4.0 (keduanya dari
        Wikimedia Commons).
\end{itemize}
\clearpage
